\section{Introduction}

The aging population in developed countries has driven demand for caregivers, making work efficiency a critical priority. Nursing homes increasingly adopt digital recording systems on smartphones and tablets to replace paper records, reducing work time and enabling data-driven analysis~\cite{okuda2024}. To complement these systems and optimize workload distribution, real-time tracking of caregiver location is essential.

Bluetooth Low Energy (BLE) beacons are widely deployed for indoor localization due to low cost, low power, and ease of installation~\cite{faragher2015}. However, BLE-based localization in real nursing environments faces two critical challenges. First, Received Signal Strength Indicator (RSSI) measurements suffer from severe noise:

\begin{equation}
\text{RSSI}(t) = P_{\text{tx}} - 10n \log_{10}(d) - \eta(t)
\label{eq:rssi_model}
\end{equation}

\noindent where $P_{\text{tx}}$ is the transmitted power, $n$ is the path-loss exponent, $d$ is the distance, and $\eta(t)$ captures multipath fading, body shadowing, and device orientation~\cite{halder2015}. Second, location data exhibits extreme class imbalance---caregivers spend disproportionate time in shared areas compared to individual rooms. Table~\ref{tab:imbalance} summarizes the five most and least frequent rooms in our dataset.

\begin{table}[h]
\centering
\caption{Top-5 and bottom-5 rooms by sample count. Nurse station outnumbers the rarest room by $407\times$.}
\label{tab:imbalance}
\begin{tabular}{lcc}
\toprule
\textbf{Room} & \textbf{Samples} & \textbf{Factor vs.\ rarest} \\
\midrule
Nurse station & 9,361 & $407\times$ \\
Cafeteria     & 5,224 & $227\times$ \\
Kitchen       & 3,102 & $135\times$ \\
Hallway       & 2,890 & $126\times$ \\
Cleaning      & 2,105 & $92\times$  \\
\addlinespace
Room 518      & 42    & $1.8\times$ \\
Room 517      & 42    & $1.8\times$ \\
Room 510      & 37    & $1.6\times$ \\
Room 505      & 44    & $1.9\times$ \\
Room 516      & 23    & $1\times$   \\
\bottomrule
\end{tabular}
\end{table}

Synthetic oversampling techniques such as SMOTE~\cite{chawla2002} are commonly used to mitigate imbalance. However, as shown in Fig.~\ref{fig:overfit}, they distort signal distributions and introduce overfitting when RSSI variance is already high---a finding consistent with recent reports from the same challenge~\cite{friedrich2024}.

\begin{figure}[h]
\centering
\fbox{\parbox{0.8\columnwidth}{\centering
\textbf{Fig.\ 1 (placeholder)} \\
Train--validation gap comparison: \\
SMOTE = 0.1459, class\_weight = 0.0997, baseline = 0.1139
}}
\caption{Overfit gap across imbalance strategies. SMOTE exhibits the largest gap, while class weighting achieves the lowest.}
\label{fig:overfit}
\end{figure}

To address these challenges, this paper proposes a non-synthetic BLE indoor localization pipeline for a 22-room nursing floor monitored by 25 beacons. Our main contributions are:

\begin{enumerate}
    \item \textbf{Savitzky-Golay filtering} for RSSI denoising---a technique well-established in spectroscopy but underexplored for BLE localization. The SG filter smooths each beacon's time series by local polynomial fitting:
    \begin{equation}
    \text{RSSI}_{\text{smooth}}[i] = \sum_{j = -m}^{m} c_j \cdot \text{RSSI}[i + j]
    \label{eq:sg}
    \end{equation}
    \noindent where $(2m + 1)$ is the window size and $c_j$ are least-squares convolution coefficients~\cite{savitzky1964}. We use $m = 2$ (window 5) and polynomial degree $p = 2$.

    \item \textbf{Sliding window temporal features} ($k = 5$) extracting mean, standard deviation, minimum, and maximum from each beacon, yielding $4 \times 25 = 100$ features per sample to capture spatial signatures across consecutive timestamps~\cite{sensors2023}.

    \item \textbf{Class-weighted Random Forest} with \texttt{balanced\_subsample} to handle extreme class imbalance ($407\times$ ratio) without generating artificial data.
\end{enumerate}

Experimental results on a real nursing facility dataset demonstrate that our best configuration (V3) achieves a Weighted F1 of 0.8351 and Macro F1 of 0.8336, with the Savitzky-Golay filter alone improving Rare-F1 by 10.3\% over raw RSSI. Table~\ref{tab:results} summarizes the performance across five explored configurations.

\begin{table}[h]
\centering
\caption{Performance comparison across five configurations. V3 (SG + class weight) is selected as the final model.}
\label{tab:results}
\begin{tabular}{lcc}
\toprule
\textbf{Configuration} & \textbf{W-F1} & \textbf{M-F1} \\
\midrule
V1: 17-class, no weight, no SG   & 0.8315 & 0.8193 \\
V2: 22-class, class weight       & 0.8246 & 0.8191 \\
\textbf{V3: 22-class, SG, class weight} & \textbf{0.8351} & \textbf{0.8336} \\
V3+4s: V3 + adaptive buffer 4s   & 0.8347 & 0.8350 \\
V4: + relabeling (cosine)        & 0.8171 & 0.7713 \\
\bottomrule
\end{tabular}
\end{table}

The remainder of this paper is organized as follows. Section~\ref{sec:related} reviews related work. Section~\ref{sec:method} describes our preprocessing pipeline and classifier. Section~\ref{sec:experiments} presents experimental setup and results. Section~\ref{sec:conclusion} concludes and discusses future directions.
